 \appendix
 \section{Some topological vector spaces}\label{appx:topspaces}

 We briefly rehearse the definition and main properties of the various $C^k$ and Sobolev spaces encountered in the text, broadly following~\cite{Baer:2015,BaerWafo:2015}, before turning to some properties of the topology of bounded convergence that are also needed. No originality is claimed for the material given here.

 \subsection[Spaces of smooth and C\^{}k functions]{Spaces of smooth and $\boldsymbol{C^k}$ functions}
 Let $M$ be a smooth manifold and let $C^k(M)$, $k\in\NN_0\cup\{\infty\}$, be the vector space of complex-valued $k$-times continuously differentiable functions on $M$. For each $k\in\NN_0$ and compact $K\subset M$, one has a seminorm
 \begin{equation*}
 	\|f\|_{K,k} = \max_{0\le r\le k} \max_{x\in K} |(\nabla^r f)(x)|_r
 \end{equation*}
 on $C^k(M)$, where $\nabla$ is an arbitrarily chosen connection on $M$ and $|\cdot|_r$ an arbitrarily chosen norm making $T^*M^{\otimes r}$ a (finite-dimensional) Banach bundle; different choices result in equivalent seminorms. The collection of seminorms $\|\cdot\|_{K,k}$ as $K$ runs over compact subsets of $M$ and $k\in\NN_0$ provides a Fr\'echet topology on $C^\infty(M)$; similarly,
 we obtain a Fr\'echet topology on~$C^k(M)$, $k\in\NN_0$, using the seminorms
 $\|\cdot\|_{K,k}$.


 If $A$ is closed, we define $C^\infty_A(M)=\{f\in C^\infty(M)\colon\supp f\subset A\}$ with the relative topology. Thus the topology is defined by the seminorms $\|\cdot\|_{K,k}$ as $K$ runs over compact subsets $K\subset A$ and $k\in\NN_0$; if $A$ is compact, it is sufficient to use the seminorms $\|\cdot\|_{A,k}$, $k\in\NN_0$. As $C^\infty_A(M)$ is closed subspace of a Fr\'echet space, it is also Fr\'echet. Defining $C^k_A(M)$ as the analogous subspace of $C^k(M)$, the topology is generated by seminorms $\|\cdot\|_{K,k}$ for compact $K\subset A$, or just the single seminorm $\|\cdot\|_{A,k}$ (which is a norm on $C^k_A(M)$) in the case that $A$ is compact.

 A \emph{support system}~\cite{Baer:2015} is a subset $\mathscr A$ of the set of all closed subsets on $M$, which is closed under finite unions and has the property that for each $A\in \mathscr A$ it holds that $(i)$
 $A\subset \textrm{int}\,(A')$ for some $A'\in\mathscr A$ and $(ii)$ if $A'$ is a closed subset of $M$ with $A'\subset A$ then $A'\in\mathscr A$. Any support system is a directed system with respect to inclusion and we write
 \begin{equation*}
 	C^\infty_{\mathscr{A}}(M)= \bigcup_{A\in\mathscr{A}} C^\infty_A(M)
 \end{equation*}
 with the locally convex inductive limit topology,
 so that a convex set $U\subset C^\infty_{\mathscr{A}}(M)$ is a neighbourhood of $0$
 if and only if $U\cap C^\infty_A(M)$ is a neighbourhood of $0$ in $C^\infty_A(M)$ for every $A\in\mathscr{A}$; because $\mathscr{A}$ is directed, one also has that
 a convex set $O$ is open if and only if $O\cap A$ is open in each $C^\infty_A(M)$.

 Examples of support systems include the set of compact sets, leading to the space of compactly supported functions $\CoinX{M}$, and the sets of (strictly) future/past/spatially-compact sets, giving rise to $C^\infty_{\sfc/\spc/\fc/\pc/\mathrm{sc}}(M)$.

 A linear map $T$ from $C^\infty_{\mathscr{A}}(M)$ to a locally convex topological vector space $Y$ is continuous if and only if all its restrictions $T_A\colon C^\infty_A(M)\to Y$ are continuous ($A\in\mathscr{A}$). In particular, for a linear map $T\colon C^\infty_{\mathscr{A}}(M)\to C^\infty_{\mathscr{B}}(M)$ to be continuous, it suffices that each restriction $T_A$ has range contained in some $C^\infty_B(M)$ for some $B\in\mathscr{B}$ (depending on $A$) and determines a continuous map $C^\infty_A(M)\to C^\infty_B(M)$, thus implying that $T_A\colon C^\infty_A(M)\to C^\infty_{\mathscr{B}}(M)$ is continuous. We record the following application:
 \begin{Lemma}%\label{lem:restr}
 	If $T$ is a continuous linear self-map of $C^\infty(M)$ such that $\supp Tf\subset K\cup \supp f$ for all $f\in C^\infty(M)$, where $K$ is a fixed compact set, then $T$ restricts to any of $C^\infty_{0/\sfc/\spc/\fc/\pc/\mathrm{sc}}(M)$ as a continuous map.
 \end{Lemma}
 \begin{proof}
 	If $\mathscr{A}$ is one of the relevant support systems then $A\cup K\in\mathscr{A}$ for each $A\in\mathscr{A}$. As the restriction $T_A$ of $T$ to $C^\infty_A(M)$ has its range in $C^\infty_{A\cup K}(M)$, and both these function spaces are subspaces of $C^\infty(M)$ with the relative topology, continuity of $T_A\colon C^\infty_A(M)\to C^\infty_{A\cup K}(M)$ follows from that of $T$. Letting $A$ vary in $\mathscr{A}$, the result is proved.
 \end{proof}

 Letting $\Omega_\alpha$ be the bundle of densities of weight $\alpha$, we can define spaces $\Gamma^\infty(\Omega_\alpha)$,
 $\Gamma_A^\infty(\Omega_\alpha)$ and $\GoinX{\Omega_\alpha}$ of smooth sections
 of $\Omega_\alpha$ in a similar way; any smooth nowhere vanishing density $\mu$ on
 $M$ provides topological isomorphisms between these spaces and their zero-weight analogues, simply by multiplication by appropriate powers of $\mu$. The space of distributions on $M$, $\DD'(M)$ is the topological dual of $\GoinX{\Omega_1}$,
 equipped with the weak-$*$ topology. In particular, $C^\infty(M)$ is canonically embedded as a subspace of $\DD'(M)$. More generally,
 $\DD'(M,\Omega_\alpha)$ is defined as the dual of $\GoinX{\Omega_{1-\alpha}}$.
 We also write $\EE'(M)$ for the topological dual of $\Gamma^\infty(\Omega_1)$, i.e.,
 the distributions of compact support.


 \subsection{Sobolev spaces}

 \subsubsection{Compact manifolds}

 {\bf The spaces $\mathbf{H^s(M)}$.}
 On a \emph{compact} manifold $M$, choose an auxiliary Riemannian metric and define $L^2(M)$ in the usual way, using the volume element induced by the metric and writing the inner product as $\langle\cdot\mid\cdot\rangle$ (linear in the second slot). Let $T=(-\triangle + I)^{1/2}$, where $\triangle$ is the Laplace--Beltrami operator, initially defined on $C^\infty(M)$ and then extended (uniquely) to a self-adjoint negative operator on $L^2(M)$.
 Its complex powers $T^s$ (as defined by functional calculus) are classical pseudodifferential operators of order $s$, $T^s\in\Psi^s(M)$~\cite{Seeley:1967} (see~\cite{Amman_etal:2004} for an axiomatically based proof) and $T^s$ is compact for $\Re s<0$ (most easily seen using the spectral properties of~$-\triangle$ on compact manifolds). The domain of $T^s$ contains $C^\infty(M)$ for all $s$.

 For $s\in\RR$, the Sobolev space $H^s(M)$ is defined as the completion of $C^\infty(M)$ with respect to the norm
 \begin{equation*}
 	\|f\|_{H^s(M)} = \| T^s f\|_{L^2(M)}
 \end{equation*}
 and, as a topological vector space, is independent of the choice of auxiliary Riemannian metric involved in its construction (of course, the specific norm $\|\cdot\|_{H^s(M)}$
 and its compatible Hilbert space inner product
 are metric-dependent). Indeed, any positive second-order elliptic operator that is essentially self-adjoint on $C^\infty(M)$ could be used in place of $-\triangle$.
 Evidently $T^s$ extends to an isometry from $H^s(M)$ to $L^2(M)$, which may be used to embed $H^s(M)$ in $\DD'(M)$ so that $u \in H^s(M)$ corresponds to the distribution $f\mapsto \ip{T^{-s}\overline{f/\nu}}{T^su}$, $f\in\GoinX{\Omega_1}$, where $\nu$ is the density corresponding to the
 volume element of the auxiliary metric used to define $L^2(M)$. Restricted to $u\in C^\infty(M)$, this embedding is consistent with the embedding of $C^\infty(M)$ in $\DD'(M)$ already mentioned.
 As $T^q$ is compact on $L^2(M)$ for $q<0$, it follows that $H^s(M)\subset H^t(M)$ for all $s>t$, with a compact inclusion map. More generally. $P\colon H^{s+m}(M)\to H^s(M)$ is continuous for any partial differential operator of order $m$ with smooth coefficients (or indeed any pseudodifferential operator $P\in \Psi^{m}(M)$), because $T^{-m}P\in \Psi^0(M)$ extends to a bounded operator on $L^2(M)$.


 {\bf Duality.} For any $s\in\RR$, suppose that $\ell\in H^s(M)'$, so there is a constant $c$ so that $|\ell(u)|\le c\|u\|_{H^s(M)}$ for all $u\in H^s(M)$. Then $|\ell(T^{-s}f)|\le c\|f\|_{L^2(M)}$ for $f\in L^2(M)$ and consequently there is $w\in L^2(M)$ so that $\ell(u) =\langle w\,|\, T^{s}u\rangle_{L^2(M)}$ for all $u\in H^s(M)$. Choose a sequence $w_n\to w$ in $L^2(M)$ with
 $w_n\in C^\infty(M)$ and note that $T^s w_n\in C^\infty(M)$ is a Cauchy sequence with respect to the
 $H^{-s}(M)$ norm, converging to some $v\in H^{-s}(M)$ for which
 $T^{-s}v=\lim_n w_n=w$. Thus
 $\ell(u)=\langle T^{-s}v\, |\,T^su\rangle$ for all $u\in H^s(M)$ and
 $|\ell(u)|\le \|v\|_{H^{-s}(M)}\|u\|_{H^s(M)}$. Noting that
 \begin{equation*}
 	\frac{\ell(T^{-s}w_n)}{\|T^{-s}w_n\|_{H^s(M)}}=\frac{\langle w\mid w_n \rangle_{L^2(M)}}{\|w_n\|_{L^2(M)}}\to \|w\|_{L^2(M)} = \|v\|_{H^{-s}(M)},
 \end{equation*}
 we see that $\|\ell\|=\|v\|_{H^{-s}(M)}$, and we have established that
 the space $H^{-s}(M)$ is anti-isomorphic to $H^s(M)'$ with the operator norm topology (i.e.,
 the strong dual) with respect to the sesquilinear pairing $\langle v,u\rangle = \langle T^{-s}v\,|\, T^s u\rangle$, $v\in H^{-s}(M)$, $u\in H^s(M)$. To be precise: every
 choice of Riemannian metric on $M$ induces an anti-isomorphism of this type, and the
 pairing depends on the choice made.

 {\bf Embedding theorems.}
 The relationship between $C^k$ and Sobolev spaces is given as follows. In one direction, the formula $-\triangle=\delta d$ (on $0$-forms) gives the estimates
 \begin{equation*}
 	0\le \ip{f}{(-\triangle)^{2r}f} = \| (-\triangle)^r f\|^2_{L^2(M)}\le C\|f\|_{M,2r}^2
 \end{equation*}
 and
 \begin{equation*}
 	0\le \ip{f}{(-\triangle)^{2r+1}f} = \| d(-\triangle)^r f\|^2_{\Lambda^1(M)}\le C\|f\|_{M,2r+1}^2 ,
 \end{equation*}
 where $\Lambda^1(M)$ is the Hilbert space of square-integrable $1$-forms (with respect to the auxiliary Riemannian metric) and $f\in C^\infty(M)$.
 From this we find
 \begin{equation*}
 	\|f\|_{H^k(M)}^2 = \ip{f}{(-\triangle+I)^k f} = \sum_{j=0}^k \binom{k}{j} \ip{f}{(-\triangle)^{j}f} \le C \|f\|_{M,k}^2
 \end{equation*}
 for all $k\in \NN_0$. Thus $C^k(M)$ is continuously embedded in $H^s(M)$ for all $s\ge k$.

 On the other hand, now let $n$ be the maximum dimension of any component of $M$.
 For $\Re s>n/2$, the integral kernel $Z_s(p,q)$ of the operator $(-\triangle+I)^{-s}=T^{-2s}$ on $L^2(M)$ is continuous, $Z_s\in C(M\times M)$ and can also be written
 in terms of the spectral decomposition of
 $(-\triangle+I)$ as
 \begin{equation*}
 	Z_s(p,q) = \sum_{r} \frac{e_r(p)\overline{e_r(q)}}{\lambda_r^{s}},
 \end{equation*}
 where $e_r\in L^2(M)$ are a basis of smooth orthonormal eigenfunctions for $-\triangle+I$ with corresponding eigenvalues $\lambda_r$ -- see~\cite{Seeley:1967}. In particular,
 \begin{equation*}
 	v_p = \sum_{r} \frac{\overline{e_r(p)}e_r}{\lambda_r^{s/2}}
 \end{equation*}
 belongs to $L^2(M)$ with $\|v_p\|^2 = Z_s(p,p)$ by the Pythagoras theorem.
 For $\Re s>n/2$, it follows that $\|e_r\|_\infty \le C_s\big|\lambda_r^{s/2}\big|$, where $C_s=\sup_{p\in M}|Z_s(p,p)|^{1/2}$. Moreover, if $f\in C^\infty(M)$,
 \begin{equation*}
 	|\langle e_r \mid f \rangle|=\lambda_r^{-k}\big|\ip{e_r}{T^k f}\big|\le \lambda^{-k}_r\|f\|_{H^k(M)},\qquad k\in\NN,
 \end{equation*}
 so $\ip{e_r}{f}$ decays faster than any inverse power of $\lambda_r$ and
 $\ip{v_p}{T^s f} = \sum_r e_r(p)\ip{e_r}{f}= f(p)$ for all (not merely almost all) $p\in M$. Consequently,
 \begin{equation*}
 	\|f\|_\infty =\sup_{p\in M} |\langle v_p \mid f \rangle|\le C_s \|f\|_{H^s(M)}, \qquad f\in C^\infty(M).
 \end{equation*}
 By density of $C^\infty(M)$ in $H^s(M)$, it follows
 that there is a continuous embedding $H^s(M)\to C(M)$ if $s>n/2$; considering $\|Pf\|_\infty$ in a~similar way for differential operators $P$, one sees that $H^s(M)$ is continuously embedded in $C^k(M)$ for all $s>k+n/2$. It follows that $\cap_{s\in\RR} H^s(M)=C^\infty(M)$.


 If $u\in \DD'(M)= \GoinX{\Omega_1}'$ is a distribution then, owing to compactness of $M$, there is a~partial differential operator $P$ so that $|u(\nu f)|\le \|Pf\|_\infty \le C_s \|Pf\|_{H^s(M)}$ for any $s>n/2$ and all $f\in C^\infty(M)$. Thus, with $t=s+\operatorname{ord} (P)$, $u(\nu f)= \big\langle \overline{f},w\big\rangle= \ip{T^t \overline{f}}{T^{-t}w}_{L^2(M)}$ for some $w\in H^{-t}(M)$, which is evidently compatible with the embedding of $H^{-t}(M)$ in $\DD'(M)$
 described earlier. In this sense,
 \begin{equation*}
 	\DD'(M) = \bigcup_{s\in\RR} H^{s}(M).
 \end{equation*}

 \subsubsection{General manifolds}

 For a general (not-necessarily compact) smooth manifold $M$ with at most finitely many components,
 we proceed as follows. If $K\subset M$ is compact and topologically regular,\footnote{Note that if $K$ is any compact subset, then the closure of the interior of $K$ is compact and topologically regular.} we may
 diffeomorphically identify $K$ with a compact subset $\hat{K}$ of a compact manifold $N$ (for example, by `doubling' a compact set that contains $K$ in its interior and has a smooth boundary~\cite{BaerWafo:2015}; for~$K$ contained in a coordinate chart one could take $N$ to be a torus). Then the Sobolev space~$H^s_K(M)$ is defined as the completion of
 $C^\infty_K(M)$ with respect to the pull back of (some choice of) the~$H^s(N)$ norm. The space $H^s_K(M)$ can be identified with a subspace of $\DD'_K(M)$, the distributions with support contained in $K$. As a topological vector space, it is independent of the choices used in its construction. However, by making such a choice one may
 endow $H^s_K(M)$ with a norm, denoted $\|\cdot\|_{H^s_K(M)}$, and indeed a compatible Hilbert space structure.
 Estimates of the form
 $\|f\|_{H^k_K(M)}\le C\|f\|_{K,k}$, $f\in C^\infty_K(M)$, and hence continuity of
 the embedding $C^k_K\to H^k_K$, carry over from the compact case for each fixed $k\in\NN_0$, as does compactness of the embedding $H^{s}_K(M)\to H^{t}_K(M)$ for $s>t$
 and the continuous embedding of $H^s_K(M)$ in $C^k_K(M)$ for $s>k+n/2$,
 where again $n$ is the maximum dimension of any component of $M$.
 Consequently one has $\cap_{s\in\RR} H^s_K(M)=C^\infty_K(M)$.

 Next, we define (abbreviating `compact and topologically regular' by c.t.r.)
 \begin{equation*}
 	H_0^s(M) = \bigcup_{\text{c.t.r. $K\subset M$}} H^s_K(M),
 \end{equation*}
 with the locally convex inductive limit topology.
 As $M$ admits countable compact exhaustions (and consequently, countable c.t.r.\ exhaustions), $H_0^s(M)$ may be realised as a countable strict inductive limit -- i.e., it is a $LF$ space. A fact to be used later is that, by~\cite[Proposition~14.6]{Treves:TVS} (see also Proposition~4 in~\cite{DieudonneSchwartz:1949}) a subset $B$ of $H_0^s(M)$ is bounded if and only if $B$ is a bounded subset of~$H^s_K(M)$ for some compact $K$. It is easily shown that
 \begin{equation*}
 	\EE'(M) = \bigcup_{s\in\RR} H_0^{s}(M).
 \end{equation*}






 Finally, the local Sobolev space $H^s_\loc(M)$ is defined as
 \begin{equation*}
 	H^s_\loc(M) = \big\{u\in\DD'(M)\colon \text{$\chi u\in H^s_0(M)$ for all $\chi\in C_0^\infty({M})$}\big\},
 \end{equation*}
 and is equipped with the Fr\'echet topology induced by the seminorms $\|\chi \cdot\|_{H^s_{K}(M)}$ as $K$ runs over compact topologically regular subsets of $M$ and $\chi$ runs over $C^\infty_K(M)$. The inclusion
 $H^s_0(M)\hookrightarrow H^s_\loc(M)$ is continuous for all $s$ and indeed we have
 \begin{equation*}
 	H^s_0(M) = H^s_\loc(M) \cap \EE'(M).
 \end{equation*}
 Thus, if $M$ is compact, $H^s_\loc(M)$ and $H^s_0(M)$ coincide as topological vector spaces. The construction above produces the same spaces as the chart-based approach taken in~\cite{Hormander3}.

 The Sobolev embedding theorems already mentioned entail the existence of continuous embeddings of $C^k(M)$ in $H^k_\loc(M)$ for all $k\in\NN_0$ and of $H^s_\loc(M)$ in $C^k(M)$ for all $s>k+n/2$, where~$n$ is the maximum dimension of any component of $M$. Consequently, $\cap_{s\in\RR} H^s_\loc(M)=C^\infty(M)$.

 \subsection{The topology of bounded convergence}\label{appx:top_bded_cvgnce}

 A general reference for the following is Tr\`eves~\cite[Chapter~32]{Treves:TVS}.

 If $E$ and $F$ are Hausdorff locally convex topological spaces then the topology of bounded convergence on $\LL(E,F)$, the space of all continuous linear maps $E\to F$, is defined by the neighbourhood base of zero, consisting of sets
 \begin{equation*}
 	\Uf(B;V) = \{T\in \LL(E,F)\colon T(B)\subset V\}
 \end{equation*}
 as $B$ runs over the bounded subsets of $E$ and $V$ runs over any neighbourhood base of zero in the topology of $F$. The notation $\LL_b(E,F)$ denotes $\LL(E,F)$ equipped with the topology of bounded convergence.
 Thus a net $T_\alpha$ converges to $0$ in $\LL_b(E,F)$
 if and only if, for every bounded $B\subset E$ and neighbourhood base set $V$,
 $T_\alpha$ eventually maps $B$ into $V$. This topology makes $\LL(E,F)$ Hausdorff and locally convex (inherited from $F$~\cite[p.~336]{Treves:TVS}). We record some basic facts.
 \begin{Lemma}\label{lem:Lb_basic}
 	Let $E$, $F$, $G$ be Hausdorff locally convex topological spaces. $(a)$ If $T\in \LL(E,F)$ and the net $S_\alpha\to S$ in $\LL_b(F,G)$ then $S_\alpha T\to ST$ in $\LL_b(E,G)$; $(b)$ if the net $T_\alpha\to T$ in $\LL_b(E,F)$ and $S\in \LL(F,G)$ then
 	$ST_\alpha\to ST$ in $\LL_b(E,G)$; $(c)$ if $T_\alpha\to T$ in $\LL_b(E,F)$ and $\Ran T_\alpha\subset \hat{F}$, where $\hat{F}$ is a topological subspace of $F$, then $T_\alpha\to T$ in $\LL_b(E,\hat{F})$; $(d)$ if $F$ is barrelled and $T_n\to T$ and $S_n\to S$ are convergent \emph{sequences} in $\LL_b(E,F)$ and $\LL_b(F,G)$, then $S_n T_n\to ST$ in~$\LL_b(E,G)$.
 \end{Lemma}
 Note that Hilbert, Banach, Fr\'echet and LF spaces are all barrelled~\cite[Chapter~33]{Treves:TVS}.
 \begin{proof}
 	$(a)$ It is enough to prove this in the case $S=0$; taking any bounded $B$ in $E$ with $0\in B$, and any $0$-neighbourhood $V$ in $G$, note that $T(B)$ is bounded in $F$ and deduce that $S_\alpha T(B)$ is eventually contained in $V$. Thus $S_\alpha T\to 0$ in $\LL_b(E,G)$. $(b)$ Without loss, suppose $T=0$ and with $B$ and $V$ as before, we note that $S^{-1}(V)$ is a $0$-neighbourhood in $F$ and again deduce that $ST_\alpha(B)$ is eventually contained in $V$, so $ST_\alpha\to 0$.
 	
 	$(c)$ Again, it is enough to treat $T=0$. Take any bounded $B\subset E$ and $0$-neighbourhood~$\hat{V}$ in~$\hat{F}$. Then $\hat{V}=\hat{F}\cap V$ for an $0$-neighbourhood~$V$ in $F$ because $\hat{F}$ carries the subspace topology. We know that, eventually, $T_\alpha(B)\subset V$, and as $\Ran T_\alpha\subset \hat{F}$, we have, eventually, that $T_\alpha(B)\subset V\cap\hat{F}= \hat{V}$. Thus $T_\alpha\to 0$ in $\LL_b(E,\hat{F})$.
 	
 	$(d)$ Suppose first that $T=0$ and let $0\in B$
 	be any bounded subset of $E$ and $V$ any $0$-neighbourhood in $G$. Then
 	the $S_n$, together with $S$, form a bounded subset in $\LL_b(F,G)$ (which, we recall, is Hausdorff). As $F$ is barrelled, they form an equicontinuous set of maps by the Banach--Steinhaus theorem~\cite[Theorem 33.1]{Treves:TVS}. Thus there exists
 	a $0$-neighbourhood $W$ in $F$ such that $S_n (W)\subset V$ for all $n$.
 	As $T_n(B)\subset W$ for all sufficiently large $n$, we have $S_nT_n(B)\subset W$ for such $n$. Thus $S_nT_n\to 0$ in this case. If $T\neq 0$, we now know
 	that $S_n(T_n-T)\to 0$, and as $S_n T\to ST$ we find $S_nT_n\to ST$ as required.
 \end{proof}

 Typical applications of Lemma~\ref{lem:Lb_basic}\,$(a)$,\,$(b)$ are to show, for instance, that a holomorphic function from a domain in $\CC$ to $\LL_b\big(H^s_0(M),H^s_K(M)\big)$ is also holomorphic as a function to $\LL_b\big(H^s_K\big)$ and
 $\LL_b(H^s_0(M))$ due to continuous embedding of $H_K^s(M)$ in $H^s_0(M)$. Part~$(d)$ has the following use.\looseness=-1
 \begin{Corollary}\label{cor:Leibniz}
 	Let $E$, $F$, $G$ be Hausdorff locally convex topological spaces, with $F$ being barrelled.
 	If $\lambda\mapsto T(\lambda)$ and $\lambda\mapsto S(\lambda)$ are functions from a domain in $\CC$ to~$\LL(E,F)$ and~$\LL(F,G)$ respectively, then $(a)$ if $T$ and $S$ are both continuous in the topology of bounded convergence $($at $\mu$$)$, so is $(ST)(\lambda)=S(\lambda)T(\lambda)$ $($at $\mu$$)$; $(b)$ if $T$ and $S$ are differentiable at $\mu$ in the topology of bounded convergence then so is $ST$, with a derivative given by the Leibniz rule
 	$(ST)'(\mu)=S'(\mu)T(\mu) + S(\mu)T'(\mu)$.
 \end{Corollary}
 \begin{proof}
 	As $\CC$ is first-countable, questions of continuity and differentiability may be reduced to sequential considerations, and the result follows using Lemma~\ref{lem:Lb_basic}\,$(d)$ and the standard proof of the Leibniz formula.
 \end{proof}

 It is also useful to have some sufficient conditions for convergence in $\LL_b(E,F)$ topology where~$E$ and $F$ are Fr\'echet or countable strict inductive limits thereof (LF spaces). If $F$ is Fr\'echet, then the bounded subsets are precisely the subsets $B$ so that $\sup_{f\in B} \rho_j(f)<\infty$ for every seminorm~$\rho_j$ defining the topology of $F$, while sets $V_{j,\epsilon}=\{f\in F\colon\rho_j(f)<\epsilon\}$, for arbitrary~$\epsilon>0$ and defining seminorm $\rho_j$, provide a basis of neighbourhoods of zero. On the other hand, if $E$ is an LF space with defining sequence $E_n$ of Fr\'echet spaces, then
 the bounded subsets~$B$ of~$E$ comprise precisely those subsets that are bounded subsets of some $E_n$~\cite[Proposition~14.6]{Treves:TVS}, while a neighbourhood base is provided by convex sets $V\subset E$ so that
 each~$V\cap E_n$ is an open neighbourhood of zero in $E_n$.
 We recall that continuous linear maps between topological vector spaces preserve boundedness~\cite[Proposition~14.2]{Treves:TVS}. The following results are used in the text.
 \begin{Lemma} \label{lem:LbFrechet}
 	Suppose $T_\alpha\in \LL_b(F)$ is a net of operators on a Fr\'echet space $F$ with defining seminorms $\rho_j$. If, for each $j$, there exists $k(j)$,
 	such that for all $\epsilon>0$, it is eventually true that $\rho_j(T_\alpha f)<\epsilon\rho_{k(j)}(f)$ for all $f\in F$, then $T_\alpha\to 0$ in $\LL_b(F)$.
 \end{Lemma}
 \begin{proof}
 	Given any bounded set $B$ and
 	neighbourhood $V_{j,\epsilon}$, we set $C=\sup_{f\in B}\rho_{k(j)}(f)$ and apply the given property to $\epsilon/C$ to find that, eventually,{\samepage
 	\begin{equation*}
 		\rho_j(T_\alpha f)<\epsilon C^{-1} \rho_{k(j)}(f) \le \epsilon, \qquad \forall f\in B.
 	\end{equation*}
 	We have shown that, eventually, $T_\alpha\in \Uf(B;V_{j,\epsilon})$; as $j$ and $\epsilon$ were arbitrary, $T_\alpha\to 0$ in $\LL_b(F)$.}
 \end{proof}

 \begin{Corollary} \label{cor:HstoCinfty}
 	Consider a net $T_\alpha\in \LL\big(C^\infty_K\big)$
 	such that each $T_\alpha$ extends to an operator in~$\LL(H^s_K)$ $($which we also denote $T_\alpha$$)$ for all $s\ge s_*$. Suppose also that $T_\alpha\to 0$ in every $\LL(H^s_K)$, $s\ge s_*$. Then $T_\alpha\to 0$ in $\LL_b\big(C^\infty_K\big)$.
 \end{Corollary}
 \begin{proof}
Any $C^\infty_K$ seminorm is bounded by an $H^s_K$ norm and vice versa, so, for any
 	$C^\infty_K$-seminorm $\|\cdot\|_{K,j}$ we may estimate
 	\begin{equation*}
 		\|T_\alpha f\|_{K,j} \le c \| T_\alpha f\|_{H^{s(j)}_K} \le c\|T_\alpha\|_{\LL(H^{s(j)}_K)} \|f\|_{H^{s(j)}_K}
 		\le c' \|T_\alpha\|_{\LL(H^{s(j)}_K)} \|f\|_{K,k(s(j))}
 	\end{equation*}
 	for suitable choices of $s(j)$ and $k(s(j))$, uniformly in $f\in C^\infty_K$ and $\alpha$.
 	As $T_\alpha\to 0$ in $\LL\big(H^{s(j)}_K\big)$, for any $\epsilon>0$ it is eventually true that
 	$\|T_\alpha f\|_{K,j} \le \epsilon \|f\|_{K,k(s(j))}$ for all $f\in C^\infty_K$. Hence $T_\alpha\to 0$ in $\LL_b\big(C^\infty_K\big)$ by Lemma~\ref{lem:LbFrechet}.
 \end{proof}

 \begin{Lemma}\label{lem:LbLF}
 	Suppose $T_\alpha$ is a net of operators on the LF space $E=\bigcup_{n\in\NN} E_n$. If for some fixed $n$, one has $\Ran(T_\alpha)\subset E_n$ and $T_\alpha|_{E_m}\to 0$ in $\LL_b(E_m,E_n)$ for all $m$, then $T_\alpha\to 0$ in both $\LL_b(E,E_n)$ and $\LL_b(E)$.
 \end{Lemma}
 \begin{proof}
 	Consider any neighbourhood $V$ of zero in the standard neighbourhood base of $E_n$ and any bounded set $B\subset E$, necessarily obeying $B\subset E_m$ for some $m$. Then $T_\alpha$ eventually maps $B$ into $V$; as $B$ and $V$ were arbitrary, $T_\alpha\to 0$ in $\LL_b(E,E_n)$. Post-composing with the continuous embedding of $E_n$ in $E$, $T_\alpha\to 0$ in $\LL_b(E)$ as well.
 \end{proof}
